\vspace*{3cm}

\begin{center}
  {\LARGE\bfseries\color{vorange!95!black}\scshape Avant-propos\par}
  \vspace{0.4cm}
  \begin{tikzpicture}
    \draw[vorange!50, line width=1pt] (-2,0) -- (2,0);
    \fill[vorange] (0,0) ++(0,2pt) -- ++(2pt,-2pt) -- ++(-2pt,-2pt) -- ++(-2pt,2pt) -- cycle;
  \end{tikzpicture}
\end{center}

\vspace{1.5cm}

\begin{center}
\begin{minipage}{0.85\textwidth}
  \setlength{\parindent}{1.5em}
  \setlength{\parskip}{0.8em}
  \large 

  {\hyphenpenalty=10000\exhyphenpenalty=10000
  Ce polycopié rassemble l'ensemble du cours de notre grand Quibel, ainsi que les travaux dirigés et certaines de leurs corrections.
  
  Il s'agit d'un travail collaboratif. Si vous souhaitez vous porter volontaire pour rédiger des corrections d'exercices ou étoffer certaines sections, faites-vous plaisir : elles seront intégrées au document.
  
  La saisie d'un tel volume de mathématiques étant sujette aux coquilles, si vous voyez une erreur, merci de nous le signaler !

  À cet effet, vous pouvez contribuer au projet sur cette repo : \url{https://github.com/mp2i-mpi-champo/mpi-maths/}.
  \par}
  \vspace{1.8cm}
  \begin{center}
    \itshape Enjoy !
  \end{center}
\end{minipage}
\end{center}

\vfill
